\chapter{مقدمه و انگیزهٔ طرح}
\label{chap:introduction}
%=====================================================================

\section{زمینه و اهمیت موضوع}
\label{sec:motivation}

\subsection{از \lr{HTTP} تا \lr{HTTPS}: چرا اصلاً باید رمز کرد؟}

در روزگاری که وب امروزی شکل می‌گرفت، بسیاری از وب‌سایت‌ها از طریق پروتکل \lr{HTTP} با کاربران خود ارتباط برقرار می‌کردند -- پروتکلی که هیچ مکانیزم رمزنگاری‌ای در آن تعبیه نشده بود. نتیجه این بود که هر کسی که در مسیر ارتباط بین مرورگر کاربر و سرور قرار می‌گرفت -- برای نمونه، فردی روی همان شبکهٔ وای‌فای عمومی، یا اپراتور یک شبکهٔ واسط -- می‌توانست محتوای رد و بدل شده را بدون هیچ مانعی بخواند و حتی، بدون آنکه هیچ‌یک از دو طرف متوجه شوند، آن را تغییر دهد. این الگوی نفوذ، که در آن مهاجم خود را نامرئی میان دو طرف ارتباط قرار می‌دهد، به \textbf{حملهٔ مرد میانی} \lr{(Man-in-the-Middle -- MITM)} معروف است. راه‌حل صنعت وب، افزودن یک لایهٔ رمزنگاری (نخست \lr{SSL} و سپس جانشین امن‌تر آن، \lr{TLS}) روی همان پروتکل بود؛ نتیجه، پروتکل \lr{HTTPS} نام گرفت -- که همان حرف \lr{S} پایانی آن، به‌روشنی معرف واژهٔ \textbf{Secure} است. اما این رمزنگاری دقیقاً چگونه کار می‌کند؟ و آیا صرفِ «رمزکردن» داده، به‌تنهایی برای امنیت واقعی کافی است؟ برای پاسخ به این پرسش‌ها، باید داستان را از پایه‌ای‌تر آغاز کرد.

\subsection{مسئلهٔ کلید مشترک: چرا رمزنگاری متقارن به‌تنهایی کافی نبود}

ساده‌ترین شکل رمزنگاری، رمزنگاری \textbf{متقارن} \lr{(Symmetric Cryptography)} است. تصور کنید \lr{Alice} می‌خواهد پیامی محرمانه برای \lr{Bob} بفرستد. ساده‌ترین راه این است که این دو، پیشاپیش بر سر یک قرارداد یا «کد» مشترک توافق کنند و \lr{Alice} هر بیت از پیام را با آن کد \lr{XOR}\LTRfootnote{عملگر منطقی جمع پیمانه‌ای دو (\lr{Exclusive OR}) که این خاصیت را دارد: اعمال دوبارهٔ آن با همان کلید، مقدار اصلی را بازمی‌گرداند. این مثال صرفاً برای درک شهودی مسئله است؛ رمزهای متقارن واقعی مانند \lr{AES} پیچیدگی به‌مراتب بیشتری دارند، هرچند اصل مسئله -- نیاز به یک کلید مشترک از پیش تعیین‌شده -- در آن‌ها نیز صدق می‌کند.} کند؛ \lr{Bob} نیز با اعمال همان عملگر \lr{XOR}، این‌بار روی دادهٔ رمزشده و با همان کد، پیام اصلی را بازمی‌یابد. این ایدهٔ ساده یک مشکل بنیادین دارد: \textbf{این کدِ مشترک را از کجا باید تهیه کرد؟} اگر \lr{Alice} و \lr{Bob} از پیش، به‌صورت حضوری و امن، این کد را با هم قرار نگذاشته باشند، تنها راهِ در دسترس این است که آن را هم از طریق همان کانال ارتباطی -- که فرض بر ناامن‌بودن آن است -- برای یکدیگر ارسال کنند. اما درست همین‌جاست که شنودگر -- که در ادبیات رمزنگاری او را \lr{Eve} می‌نامند و فرض بر آن است که میان \lr{Alice} و \lr{Bob} نشسته است -- می‌تواند همان کد را در لحظهٔ تبادل، شنود و ذخیره کند. از آن پس، \lr{Eve} دقیقاً همان کلیدی را در اختیار دارد که \lr{Bob} دارد، و کل زحمت رمزنگاری بی‌اثر می‌شود. این معضل در ادبیات فنی \textbf{مسئلهٔ توزیع کلید} \lr{(Key Distribution Problem)} نامیده می‌شود.

\subsection{راه‌حل: کلید را دوتا کنید -- تولد رمزنگاری کلید عمومی}

راه‌حلی که ریاضی‌دانان و مهندسان دهه‌ی \lr{1970} میلادی برای این بن‌بست یافتند، به‌جای تلاش برای امن‌کردن توزیع یک کلید مشترک، خودِ پیش‌فرض مسئله را به چالش کشید: \textbf{چرا اصلاً کلیدِ قفل‌کردن باید همان کلیدِ بازکردن باشد؟} دیفی و هلمن در سال \lr{1976} \Cite{diffiehellman1976} این ایده را به‌صورت مفهومی مطرح کردند: هر گیرنده (برای نمونه \lr{Bob}) یک جفت کلید ریاضی مرتبط اما نامتقارن تولید می‌کند -- یک \textbf{کلید عمومی} \lr{(Public Key)} که آن را آشکارا و بدون هیچ نگرانی برای همه (از جمله هر فرستنده‌ای و حتی خودِ \lr{Eve}) منتشر می‌کند، و یک \textbf{کلید خصوصی} \lr{(Private Key)} که تنها نزد خود نگه می‌دارد و هرگز آن را با کسی به اشتراک نمی‌گذارد. هر کس می‌تواند با کلید عمومیِ منتشرشدهٔ \lr{Bob} پیامی را رمز کند، اما تنها کسی که کلید خصوصی متناظر را در اختیار دارد -- یعنی خودِ \lr{Bob} -- قادر به رمزگشایی آن است. به این ترتیب، دیگر هیچ کلید حساسی از طریق کانال ارتباطی رد و بدل نمی‌شود، و مسئلهٔ توزیع کلید، دست‌کم در ظاهر، حل‌شده به نظر می‌رسید.

\subsection{\lr{RSA}: پیاده‌سازی ریاضی یک ایدهٔ زیبا}

ایدهٔ دیفی و هلمن، مفهومی بود؛ پیاده‌سازی ریاضی و عملی آن دو سال بعد، در سال \lr{1978}، توسط سه پژوهشگر دانشگاه \lr{MIT} -- ریوست\LTRfootnote{Ron Rivest}، شامیر\LTRfootnote{Adi Shamir} و ادلمن\LTRfootnote{Leonard Adleman} -- ارائه شد؛ نام الگوریتم \lr{RSA} نیز برگرفته از حرف نخست نام خانوادگی همین سه پژوهشگر است \Cite{rsa1978}. منطق \lr{RSA} چنین است: دو عدد اول بسیار بزرگ (هرکدام با صدها رقم) به‌طور تصادفی انتخاب و در یکدیگر ضرب می‌شوند تا عددی به‌مراتب بزرگ‌تر -- که بخشی از کلید عمومی را تشکیل می‌دهد -- به دست آید. برای یک رایانه، ضرب‌کردن این دو عدد اول کاری آنی و بی‌دردسر است؛ اما مسیر معکوس -- یعنی داشتن تنها همان عدد بزرگِ حاصل‌ضرب، و تلاش برای یافتن دو عامل اول تشکیل‌دهندهٔ آن (که برای بازسازی کلید خصوصی لازم است) -- از نظر محاسباتی، مسئله‌ای به‌شدت دشوار است. برای یک پیمانهٔ \lr{2048} بیتی (معادل تقریباً \lr{617} رقم اعشاری)، سطح امنیت کلاسیک این طرح در حدود \lr{112} بیت برآورد می‌شود \Cite{nistsp800-57}. برای درکی شهودی از این دشواری: بهترین الگوریتم کلاسیک شناخته‌شده برای این نوع تجزیه (الک میدان عددی تعمیم‌یافته \lr{(GNFS)}) حتی برای اعدادی به‌مراتب کوچک‌تر از این اندازه، به زمانی در حدود سن جهان هستی نیاز دارد \Cite{bub2006}. به بیان دیگر، امنیت \lr{RSA} بر یک اثبات ریاضی قطعی استوار نیست، بلکه صرفاً بر این باور عملی متکی است که \textbf{هیچ الگوریتم کلاسیک کارآمدی برای تجزیهٔ اعداد صحیح بزرگ در زمان چندجمله‌ای وجود ندارد.}

\subsection{الگوریتم شور: وقتی این باور فرو ریخت}

این باورِ چند دهه‌ای، در سال \lr{1994} به لرزه درآمد. پیتر شور\LTRfootnote{Peter W. Shor}، ریاضی‌دان آزمایشگاه بل، نشان داد که با بهره‌گیری از یک خاصیت صرفاً کوانتومی -- توانایی رایانهٔ کوانتومی در اجرای \textbf{تبدیل فوریهٔ کوانتومی} \lr{(Quantum Fourier Transform)} برای یافتن دورهٔ تناوب یک تابع، به‌طور هم‌زمان روی برهم‌نهی تمام مقادیر ممکن ورودی -- می‌توان مسئلهٔ تجزیهٔ اعداد صحیح بزرگ (و به همان ترتیب، مسئلهٔ لگاریتم گسسته که پایهٔ امنیتی \lr{ECC} است) را در زمان چندجمله‌ای حل کرد \Cite{shor1997}؛ این نتیجه، دو سال بعد به‌طور رسمی در مجلهٔ \lr{SIAM Journal on Computing} منتشر شد. برآوردهای مهندسی به‌روزتر نشان می‌دهند که یک رایانهٔ کوانتومی خطاپذیر با حدود \lr{20} میلیون کیوبیت فیزیکی، می‌تواند همان مسئله‌ای را که رایانه‌های کلاسیک برای حل آن به زمانی در حدود سن جهان هستی نیاز داشتند، در حدود \lr{8} ساعت حل کند \Cite{gidneyekera2019}. به بیان ساده، الگوریتم شور -- به شرط ساخته‌شدن چنین رایانه‌ای -- پایهٔ ریاضیِ امنیتِ \lr{RSA} و \lr{ECC} را به‌طور کامل بی‌اثر می‌سازد.

\subsection{تهدید «برداشت اکنون، رمزگشایی بعداً»}

اگرچه ساخت یک رایانهٔ کوانتومی «بحرانی از منظر رمزنگاری» \lr{(Cryptographically Relevant Quantum Computer -- CRQC)} همچنان با چالش‌های مهندسی جدی روبه‌روست، افق زمانی برآوردشده توسط بسیاری از نهادهای امنیتی و استانداردسازی (از جمله \lr{NIST} و \lr{NSA}\LTRfootnote{National Security Agency (United States)}) برای ظهور چنین سامانه‌ای، در بازهٔ $10$ تا $20$ سال آینده است. با این حال، تهدید واقعی بسیار زودتر آغاز می‌شود؛ زیرا در الگوی حملهٔ \textbf{«برداشت اکنون، رمزگشایی بعداً»} \lr{(Harvest Now, Decrypt Later -- HNDL)}، مهاجم می‌تواند داده‌های رمزشدهٔ حساس (نظیر تراکنش‌های بانکی، مکاتبات دولتی و اطلاعات هویتی) را همین امروز -- درست مانند همان شنودگری که در مثال \lr{XOR} بالا در کمین نشسته بود -- رهگیری و ذخیره کند تا در آیندهٔ نه‌چندان دور، پس از دسترسی به رایانهٔ کوانتومی، آن‌ها را رمزگشایی نماید. برای داده‌هایی که ارزش محرمانگی‌شان از افق زمانی ظهور \lr{CRQC} فراتر می‌رود -- نظیر اسناد نظامی، مکالمات محرمانهٔ سیاسی و اطلاعات هویتی درازمدت -- این یعنی افشای اطلاعات، حتی اگر امروز غیرممکن به نظر برسد، صرفاً به تعویق افتاده است. این واقعیت، ضرورت اقدام پیشدستانه را دوچندان می‌کند.

\subsection{دو راهبرد مقابله: سخت‌ترکردن الگوریتم، یا تکیه بر فیزیک}

جامعهٔ رمزنگاری و فیزیک، در پاسخ به این تهدید، از دو مسیر کاملاً متفاوت وارد عمل شدند:

\begin{itemize}
\item \textbf{راهبرد نخست، رویکردی الگوریتمی است:} به‌جای تکیه بر مسائلی مانند تجزیهٔ اعداد صحیح -- که الگوریتم شور آن‌ها را می‌شکند -- از خانواده‌ای دیگر از مسائل ریاضی دشوار استفاده می‌شود که، تا آنجا که دانش امروز نشان می‌دهد، حتی رایانه‌های کوانتومی نیز راه‌حل کارآمدی برای آن‌ها ندارند. این راهبرد \textbf{رمزنگاری پساکوانتومی} \lr{(Post-Quantum Cryptography -- PQC)} نام گرفته و موضوع فصل \ref{chap:pqc} است.
\item \textbf{راهبرد دوم، رویکردی فیزیکی است:} به‌جای تکیه بر یک فرض محاسباتی که ممکن است روزی -- همان‌گونه که برای \lr{RSA} رخ داد -- شکسته شود، امنیت کلید مستقیماً از قوانین بنیادین فیزیک کوانتوم، نظیر اثر مشاهده‌گر \lr{(Observer Effect)} و قضیهٔ عدم شبیه‌سازی کوانتومی \lr{(No-Cloning Theorem)}، تأمین می‌شود؛ هرگونه تلاش \lr{Eve} برای شنود کانال کوانتومی، برخلاف شنود کانال کلاسیک در مثال \lr{XOR}، اختلالی قابل‌کشف در آن پدید می‌آورد و \lr{Alice} و \lr{Bob} می‌توانند وجود شنودگر را تشخیص داده و پیش از افشای کلید، اقدام پیشگیرانه انجام دهند. این راهبرد \textbf{توزیع کلید کوانتومی} \lr{(Quantum Key Distribution -- QKD)} نام دارد و موضوع فصل‌های \ref{chap:qkdtheory} و \ref{chap:qkdarchitecture} است.
\end{itemize}

این دو راهبرد، آن‌گونه که در فصل \ref{chap:hybrid} با جزئیات بیشتری تشریح خواهد شد، رقیب یکدیگر نیستند؛ این طرح هر دو را به‌صورت مکمل، در قالب یک معماری واحدِ \textbf{دفاع در عمق} \lr{(Defense in Depth)} به کار می‌گیرد.

\section{چرا اکنون؟}

سه روند هم‌زمان، این پروژه را از یک تمرین آکادمیک به یک اولویت راهبردی تبدیل می‌کند:

\begin{enumerate}
\item \textbf{نهایی‌شدن استانداردهای جهانی:} در تاریخ \lr{13} آگوست \lr{2024}، مؤسسهٔ ملی استاندارد و فناوری آمریکا \lr{(NIST)} پس از فرآیندی هشت‌ساله، سه استاندارد نهایی رمزنگاری پساکوانتومی -- \lr{ML-KEM (FIPS 203)}، \lr{ML-DSA (FIPS 204)} و \lr{SLH-DSA (FIPS 205)} -- را منتشر کرد \Cite{nist2024release}، که مسیر مهاجرت جهانی را از حالت نظری به عملیاتی تبدیل کرده است.
\item \textbf{بلوغ صنعتی \lr{QKD}:} فناوری \lr{QKD} از آزمایشگاه به مقیاس ملی و بین‌قاره‌ای رسیده است؛ نمونهٔ برجستهٔ آن کریدور $2000$ کیلومتری پکن--شانگهای در چین \Cite{mdpi2023qkdn} و پیوند اخیر $303$ کیلومتری در سوئد است \Cite{arxiv2606sweden}.
\item \textbf{ورود جدی بخش بانکی:} بانک‌های پیشرو جهان نظیر \lr{HSBC} و \lr{JPMorgan Chase} دیگر در فاز آزمایشگاهی نیستند و مستقیماً معاملات واقعی و شبیه‌سازی‌شده را روی کانال‌های امن کوانتومی اجرا می‌کنند \Cite{techinformed2025banks}.
\end{enumerate}

\section{اهداف طرح}

اهداف اصلی این پروپوزال به شرح زیر است:

\begin{itemize}
\item ارائهٔ تصویری جامع و فنی از فناوری \lr{QKD} و \lr{PQC} برای تصمیم‌گیران سیاست‌گذاری و فنی کشور؛
\item طراحی معماری یک شبکهٔ آزمایشی \lr{(Pilot)} کلان‌شهری در تهران با تمرکز بر اتصال بانک مرکزی، مراکز دادهٔ بانکی و نهادهای حساس دولتی؛
\item ارائهٔ برآورد هزینهٔ تجهیزات، زیرساخت فیبر و نگهداشت بر مبنای قیمت‌های واقعی بازار جهانی؛
\item تدوین نقشهٔ راه اجرایی چندمرحله‌ای همراه با تحلیل ریسک امنیتی و پیشنهادهای کاهش آسیب‌پذیری؛
\item زمینه‌سازی برای انتقال دانش فنی، بومی‌سازی تدریجی تجهیزات و تربیت نیروی متخصص داخلی.
\end{itemize}

%=====================================================================
